\subsection{\texorpdfstring{映射 \(\Theta_{a}\) 的性质}{映射 \textbackslash Theta\_\{a\} 的性质}}

回顾（VAR, R2, p.~46 与 p.~47, 11.1.3, d)）：若 K 是 C 的一个紧子集，则存在由 K 的紧邻域组成的基本系，其中的邻域都是 C 的块（也就是说，对每个 \(x \in V\)，存在 C 在 x 处的一个卡 \((\mathrm{U}, \varphi, \mathbf{C})\)，使得 \(\varphi(\mathrm{U} \cap \mathrm{K})\) 是 C 的某个闭半空间的开子集）。我们用 \(\partial V\) 表示赋予由 C 的定向所诱导的定向后的块 V 的边界 (VAR, R2, p.~47)，并用 dz 表示内射 \(\partial V \to C\) 的微分。

命题 1. --- 设 a 是 A 的一个元素，U 是 Sp(a) 的一个开邻域，并设 \(f \in \mathcal{O}(\mathrm{U};\mathrm{A})\)。设 V 是 Sp(a) 的一个包含于 U 的紧邻域，且 V 是 C 的一个块。

则映射 \(z \mapsto f(z)(z - a)^{-1}\) 在 \(\partial\mathrm{V}\) 上连续；微分形式 \(f(z)(z - a)^{-1}dz\) 在 \(\partial\mathrm{V}\) 上可积，并且有

\[
\Theta_ {a} (f) = \frac {1}{2 i \pi} \int_ {\partial \mathrm{V}} f (z) (z - a) ^ {- 1} d z.
\]

设 \(h\) 是从 \(\mathbf{C}\) 到 \(\mathbf{C}\) 中的无穷次可微映射，它在 \(\mathrm{Sp}(a)\) 的一个邻域中等于 1，并且其紧支集包含于 \(V\) 的内部 (I, p.~52, 引理 1)。设 \(u\) 是从 \(\mathbf{C}\) 到 A 中的映射，使得二元组 \((h, u)\) 适配于 \(a\)（参见 I, p.~53, 例 3）。相伴的微分形式是 \(\omega = du \wedge dz\)。由此得 \(f\omega = fdu \wedge dz = d(fu dz)\)，因为 \(f\) 是全纯的。此外，\(u(z) = (z - a)^{-1}\) 在 \(V\) 的边界上成立。再者，微分形式 \(fu dz\) 在 U 上是 \(C^1\) 类的。于是，

\[
2 i \pi \Theta_ {a} (f) = \int_ {\mathrm{V}} d (f u d z) = \int_ {\partial \mathrm{V}} f u d z = \int_ {\partial \mathrm{V}} f (z) (z - a) ^ {- 1} d z
\]

这是由公式 (5) 以及关于块 V 的 Stokes 公式 (VAR, R2, p.~47, 11.2.3) 得到的。

\[
\text { COROLLARY. } - \text {   Let   } a \in \mathrm{A}. \text {   One has   } \Theta_ {a} (1) = 1.
\]

设 \(\mathrm{R} > \varrho(a)\) 是一个实数。设 V 是以 0 为中心、以 R 为半径的闭圆盘，于是 \(\mathrm{Sp}(a) \subset \mathring{\mathrm{V}}\)。它是 C 的一个块，其边界 \(\partial\mathrm{V}\) 是以 0 为中心、以 R 为半径的圆周。对 \(z \in \mathbf{C} - \mathring{\mathrm{V}}\)，我们有公式 \((z - a)^{-1} = z^{-1}(1 - z^{-1}a)^{-1} = \sum_{j=1}^{+\infty} z^{-j} a^{j-1}\)。该级数对 \(z \in \partial V\) 一致收敛。因此我们有

\[
\Theta_ {a} (1) = \frac {1}{2 i \pi} \sum_ {j = 1} ^ {+ \infty} a ^ {j - 1} \int_ {\partial \mathrm{V}} z ^ {- j} d z = 1
\]

因为

\[
\int_ {\partial \mathrm{V}} z ^ {j} d z = 0
\]

对每个整数 \(j \neq -1\) 成立，且

\[
\int_ {\partial \mathrm{V}} z ^ {- 1} d z = 2 i \pi
\]

(VAR, R2, p.~44, 10.4.5 及 p.~47, 11.2.1, 例)。

引理 8. --- 设 U 是 \(C^{n}\) 的一个开子集。设 \(\omega_{1}\) 是 U 上的一个 n 次连续微分形式，具有紧支集且取值于 A（相应地，设 \(\omega_{2}\) 是 C 上的一个 2 次连续微分形式，具有紧支集且取值于 C）。我们用 \(\pi_{1}\) 与 \(\pi_{2}\) 表示 U × C 到 U 与 C 上的典范投影。U × C 上的微分形式 \(\pi_{1}^{*}\omega_{1} \wedge \pi_{2}^{*}\omega_{2}\) 连续、具有紧支集且取值于 A。我们有

\[
\int_ {\mathrm{U} \times \mathbf {C}} \pi_ {1} ^ {*} \omega_ {1} \wedge \pi_ {2} ^ {*} \omega_ {2} = \left(\int_ {\mathbf {C}} \omega_ {2}\right) \left(\int_ {\mathrm{U}} \omega_ {1}\right).
\]

设 \(\mu_{n}\) 是 \(C^{n}\) 上的 Lebesgue 测度，\(\mu_{1}\) 是 C 上的 Lebesgue 测度。存在 \(\psi_{1} \in \mathcal{K}(U;A)\) 与 \(\psi_{2} \in \mathcal{K}(C)\)，使得与 \(\omega_{1}\) 相伴的向量测度等于 \(\psi_{1} \cdot \mu_{n}\)，而与 \(\omega_{2}\) 相伴的向量测度等于 \(\psi_{2} \cdot \mu_{1}\)。与微分形式 \(\pi_{1}^{*}\omega_{1} \wedge \pi_{2}^{*}\omega_{2}\) 相伴的向量测度是 \((\psi_{1} \otimes \psi_{2}) \cdot \mu_{n} \otimes \mu_{1}\)。

设 \(\ell\) 是 A 上的一个连续线性形式。由 INT, VI, §2, 第 2 小节, 定义 2 以及乘积测度的定义 (INT, III, §4, 第 1 小节, 定义 1) 得

\[
\begin{array}{r l} & {\ell \Big (\int_ {\mathrm{U} \times \mathbf {C}} \pi_ {1} ^ {*} \omega_ {1} \wedge \pi_ {2} ^ {*} \omega_ {2} \Big) = \int_ {\mathrm{U} \times \mathbf {C}} \ell \circ (\psi_ {1} \otimes \psi_ {2}) \mu_ {n} \otimes \mu_ {1}} \\ & {\qquad = \int_ {\mathrm{U} \times \mathbf {C}} \psi_ {2} (z) \ell (\psi_ {1} (x)) d \mu_ {n} (x) d \mu_ {1} (z)} \\ & {\qquad = \Big (\int_ {\mathbf {C}} \psi_ {2} (z) d \mu_ {1} (z) \Big) \Big (\int_ {\mathrm{U}} \ell (\psi_ {1} (x)) d \mu_ {n} (x) \Big)} \\ & {\qquad = \Big (\int_ {\mathbf {C}} \psi_ {2} \mu_ {1} \Big) \ell \Big (\int_ {\mathrm{U}} \psi_ {1} \mu_ {n} \Big),} \end{array}
\]

由此即得结论 (INT VI, 同前)。

引理 9. --- 设 \(\boldsymbol{a} = (a_1, \ldots, a_n) \in \mathbb{A}^n\)。设 \(p \in \mathbf{N}\)，\(\boldsymbol{a}' = (a_{n+1}, \ldots, a_{n+p}) \in \mathbb{A}^p\)。则我们有 \(\Theta_{(\boldsymbol{a}, \boldsymbol{a}')} \circ \pi_{n,n+p}^* = \Theta_{\boldsymbol{a}}\)。特别地，我们有 \(\Theta_{\boldsymbol{a}}(1) = 1\)。

由于 \(\pi_{n,n+p} = \pi_{n,n+1} \circ \cdots \circ \pi_{n+p-1,n+p}\)，只需对 \(p = 1\) 的情形证明第一个断言，我们现在就作此假设。我们简记 \(\pi = \pi_{n,n+1}\)。于是只需证明：对 \(\mathrm{Sp}^n(\boldsymbol{a})\) 的每个开邻域 U 以及每个函数 \(f \in \mathcal{O}(\mathrm{U};\mathrm{A})\)，有 \(\Theta_{(\boldsymbol{a},a_{n+1})}(f \circ \pi) = \Theta_{\boldsymbol{a}}(f)\)。记 \(g = f \circ \pi\)。设 \(h\)（相应地，\(h'\)）是从 \(\mathbf{C}^n\) 到 \(\mathbf{C}\) 中（相应地，从 \(\mathbf{C}\) 到 \(\mathbf{C}\) 中）的无穷次可微映射，它在 \(\mathrm{Sp}^n(\boldsymbol{a})\)（相应地，\(\mathrm{Sp}(\boldsymbol{a}')\)）的一个邻域中等于 1，且具有包含于 U（相应地，\(\mathbf{C}\)）中的紧支集。存在映射 ( \(u_1, \ldots, u_n\) )：从 \(\mathbf{C}^n\) 到 A 中、无穷次可微，使得序列 ( \(h, u_1, \ldots, u_n\) ) 适配于 \(\boldsymbol{a}\) (I, p.~54, 引理 3)；又存在无穷次可微映射 \(u_{n+1}\)：从 \(\mathbf{C}\) 到 A 中，使得二元组 ( \(h', u_{n+1} \) ) 适配于 \(a_{n+1}\)（同前）。

对 \(z \in C^{n}\) 与 \(z_{n+1} \in C\)，令 \(h''(z, z_{n+1}) = h(z)h'(z_{n+1})\)，\(u_{n+1}''(z, z_{n+1}) = h(z)u_{n+1}(z_{n+1})\)。函数 \(h''\) 与 \(u_{n+1}''\) 在 \(C^{n+1}\) 中无穷次可微。函数 \(h''\) 在 \(\mathrm{Sp}^{n+1}(a, a_{n+1})\) 的一个邻域中等于 1，且具有包含于 \(U \times C\) 的紧支集。对每个 \(\boldsymbol{w} = (z, z_{n+1}) \in \mathbf{C}^{n+1}\)，有

\[
\begin{array}{c} (z _ {1} - a _ {1}) (u _ {1} \circ \pi) (\boldsymbol {w}) + \dots + (z _ {n} - a _ {n}) (u _ {n} \circ \pi) (\boldsymbol {w}) \\ + (z _ {n + 1} - a _ {n + 1}) u _ {n + 1} ^ {\prime \prime} (\boldsymbol {w}) = 1 - h (\boldsymbol {z}) + h (\boldsymbol {z}) (1 - h ^ {\prime} (z _ {n + 1})) \\ = 1 - h ^ {\prime \prime} (\boldsymbol {w}), \end{array}
\]

这表明序列 \((h''_{1}, u_{1} \circ \pi, \ldots, u_{n} \circ \pi, u_{n+1}')\) 适配于 \((\boldsymbol{a}, a_{n+1})\)。设 \(\omega\) 是相伴的微分形式。

微分形式 \(du_{1} \wedge dz_{1} \wedge \cdots \wedge du_{n} \wedge dz_{n} \wedge dh\)（它定义在 \(C^{n}\) 上）的次数为 \(2n + 1\)，故为零。因此，

\[
\begin{array}{l} \omega = d (u _ {1} \circ \pi) \wedge d z _ {1} \wedge \dots \wedge d (u _ {n} \circ \pi) \wedge d z _ {n} \wedge d u _ {n + 1} ^ {\prime \prime} \wedge d z _ {n + 1} \\ = (h \circ \pi) d (u _ {1} \circ \pi) \wedge d z _ {1} \wedge \dots \wedge d (u _ {n} \circ \pi) \wedge d z _ {n} \wedge d u _ {n + 1} \wedge d z _ {n + 1}. \end{array}
\]

由于 \(g = f \circ \pi\)，公式 (5) 与引理 8 蕴涵

\[
\begin{array}{c} \Theta_ {\boldsymbol {a}, \boldsymbol {a} ^ {\prime}} (g) = \frac {(n + 1) !}{(2 i \pi) ^ {n + 1}} \int_ {\mathrm{U} \times \mathbf {C}} g \omega = \frac {(n + 1) !}{(2 i \pi) ^ {n + 1}} \Big (\int_ {\mathbf {C}} d u _ {n + 1} \wedge d z _ {n + 1} \Big) \\ \qquad \qquad \times \Big (\int_ {\mathrm{U}} f h   d u _ {1} \wedge d z _ {1} \wedge \dots \wedge d u _ {n} \wedge d z _ {n} \Big). \end{array}
\]

一方面，我们有

\[
\int_ {\mathbf {C}} d u _ {n + 1} \wedge d z _ {n + 1} = 2 i \pi \Theta_ {a _ {n + 1}} ^ {\mathbf {C}} (1) = 2 i \pi \cdot 1
\]

（由推论 5）。另一方面，I, p.~54, 引理 4 的 c) 以及闭形式的积分为零这一事实 (VAR, R2, p.~48, 11.2.4) 蕴涵

\[
\begin{array}{c} (n + 1) \int_ {\mathrm{U}} f h d u _ {1} \wedge d z _ {1} \wedge \dots \wedge d u _ {n} \wedge d z _ {n} = \\ \int_ {\mathrm{U}} f   d u _ {1} \wedge d z _ {1} \wedge \dots \wedge d u _ {n} \wedge d z _ {n} = \frac {(2 i \pi) ^ {n}}{n !} \Theta_ {\boldsymbol {a}} (f). \end{array}
\]

这样就得到

\[
\Theta_ {\boldsymbol {a}, \boldsymbol {a} ^ {\prime}} (g) = \frac {(n + 1) !}{(2 i \pi) ^ {n + 1}} \times \frac {(2 i \pi) ^ {n}}{(n + 1) !} \Theta_ {\boldsymbol {a}} (f) \times 2 i \pi = \Theta_ {\boldsymbol {a}} (f).
\]

最后，公式 \(\Theta_{\boldsymbol{a}}(1)=1\) 由前述结果以及命题 1 的推论得出。

引理 10. --- 设 \(\boldsymbol{a} \in \mathrm{A}^n\)。设 \(g\) 是 \(\mathbf{C}^n\) 上的一个以 A 为系数的多项式函数，并设 \(f \in \mathcal{O}(\mathrm{Sp}^n(\boldsymbol{a});\mathrm{A})\)。我们有 \(\Theta_{\boldsymbol{a}}(gf) = g(\boldsymbol{a})\Theta_{\boldsymbol{a}}(f)\)。特别地，我们有 \(\Theta_{\boldsymbol{a}}(g) = g(\boldsymbol{a})\)。

由引理 9，只需证明第一个断言。

我们用 \(z_{1},\ldots,z_{n}\) 表示 \(C^{n}\) 上的坐标函数。由于映射 \(\Theta_{a}\) 是 A-线性的，只需对 \(g=z_{1}^{e_{1}}\cdots z_{n}^{e_{n}}\)（其中 \((e_{1},\ldots,e_{n})\in\mathbf{N}^{n}\)）的情形证明引理的断言。对 \(e_{1}+\cdots+e_{n}\) 作归纳，可化归为存在整数 i 使得 \(1\leqslant i\leqslant n\) 且 \(g=z_{i}\) 的情形。

设 U 是 \(\mathrm{Sp}^{n}(\boldsymbol{a})\) 的一个开邻域。设 \((h,u_{1},\ldots,u_{n})\) 是适配于 \(\boldsymbol{a}\) 的序列，且 h 的支集包含于 U (I, p.~52, 引理 1 及 I, p.~54, 引理 3)，并设 \(\omega\) 是相伴的微分形式。由 I, p.~54, 引理 4, a)，存在一个微分形式 \(\beta\)，使得

\[
(z _ {i} - a _ {i}) \omega = d (h \beta \wedge d z _ {1} \wedge \dots \wedge d z _ {n}).
\]

因此，对每个函数 \(f \in \mathcal{O}(\mathrm{U};\mathrm{A})\)，有

\[
\left(z _ {i} - a _ {i}\right) f \omega = f d \left(h \beta \wedge d z _ {1} \wedge \dots \wedge d z _ {n}\right) = d \left(f h \beta \wedge d z _ {1} \wedge \dots \wedge d z _ {n}\right)
\]

这是因为 f 全纯，从而 \(df \wedge dz_{1} \wedge \cdots \wedge dz_{n} = 0\)。应用 Stokes 公式 (VAR, R2, p.~48, 11.2.4)，我们得到 \(\int_{\mathrm{U}}(z_{i} - a_{i})f\omega = 0\)，由此得

\[
\Theta_ {\boldsymbol {a}} (z _ {i} f) = \frac {n !}{(2 i \pi) ^ {n}} \int_ {\mathrm{U}} z _ {i} f \omega = \frac {n !}{(2 i \pi) ^ {n}} \int_ {\mathrm{U}} a _ {i} f \omega = a _ {i} \Theta_ {\boldsymbol {a}} (f)
\]

（由公式 (5)）。结论得证。

命题 2. — 设 \(\varrho_{1},\ldots,\varrho_{n}\) 是 \(>0\) 的实数，并设 \(\mathrm{U}\subset\mathbf{C}^{n}\) 是以 0 为中心、以 \(\varrho_{i}\) 为半径的开圆盘的乘积多圆柱。设

\[
\sum_ {(k _ {1}, \dots , k _ {n}) \in \mathbf {N} ^ {n}} c (k _ {1}, \dots , k _ {n}) \mathrm{X} _ {1} ^ {k _ {1}} \dots \mathrm{X} _ {n} ^ {k _ {n}} \in \mathrm{A} [ [ \mathrm{X} _ {1}, \dots , \mathrm{X} _ {n} ] ]
\]

是一个以 A 为系数的形式级数。假设该级数在 U 中收敛，并用 f 表示作为其和的 U 中的全纯函数。

设 \(\boldsymbol{a}=(a_{1},\ldots,a_{n})\in\mathrm{A}^{n}\) 满足 \(\varrho(a_{i})<\varrho_{i}\)（对 \(1\leqslant i\leqslant n\)）。则 \(\mathrm{Sp}^{n}(\boldsymbol{a})\subset\mathrm{U}\)，A 的元素所成的族 \((c(k_{1},\ldots,k_{n})a_{1}^{k_{1}}\cdots a_{n}^{k_{n}})\) 绝对可和，且

\[
\Theta_ {\boldsymbol {a}} (f) = \sum_ {(k _ {1}, \dots , k _ {n}) \in \mathbf {N} ^ {n}} c (k _ {1}, \dots , k _ {n}) a _ {1} ^ {k _ {1}} \dots a _ {n} ^ {k _ {n}}.
\]

对 A 的每个特征标 \(\chi\) 以及每个整数 \(i\)（满足 \(1 \leqslant i \leqslant n\)），有 \(|\chi(a_i)| \leqslant \varrho(a_i) < \varrho_i\)，于是由联合谱的定义得 \(\mathrm{Sp}^n(\boldsymbol{a}) \subset \mathrm{U}\)。设 \(z_1, \ldots, z_n\) 是 \(\mathbf{C}^n\) 的坐标函数在 U 上的限制。则族 \((c(k_1, \ldots, k_n)z_1^{k_1} \cdots z_n^{k_n})\) 在 \(\mathcal{O}(\mathrm{U}; \mathrm{A})\) 中可和，且和为 \(f\)。顾及引理 10 以及映射 \(\Theta_{\boldsymbol{a}}^{\mathrm{U}}\) 的连续性，族 \((c(k_1, \ldots, k_n)a_1^{k_1} \ldots a_n^{k_n})\) 因而在 A 中可和，且和为 \(\Theta_{\boldsymbol{a}}(f)\)。对 \(1 \leqslant i \leqslant n\)，设 \(\lambda_i\) 是满足 \(\varrho(a_i) < \lambda_i < \varrho_i\) 的实数。存在 \(\mathrm{M}_i < +\infty\)，使得 \(\|a_i^k\| \leqslant \mathrm{M}_i\lambda_i^k\)（对每个整数 \(k \geqslant 0\)）。于是我们有

\[
\sum_ {(k _ {1}, \dots , k _ {n}) \in \mathbf {N} ^ {n}} \| c (k _ {1}, \dots , k _ {n}) \| \| a _ {1} ^ {k _ {1}} \dots a _ {n} ^ {k _ {n}} \| \leqslant
\]

\[
\mathrm{M} _ {1} \dots \mathrm{M} _ {n} \sum_ {(k _ {1}, \dots , k _ {n}) \in \mathbf {N} ^ {n}} \| c (k _ {1}, \dots , k _ {n}) \| \lambda_ {1} ^ {k _ {1}} \dots \lambda_ {n} ^ {k _ {n}}
\]

它依假设是有限的，故族 \((c(k_{1},\ldots,k_{n})a_{1}^{k_{1}}\cdots a_{n}^{k_{n}})\) 绝对可和。

推论. --- 假设 A 非零。设 \(a \in C^{n} \subset A^{n}\)。我们有 \(\operatorname{Sp}_{\mathrm{A}}^{n}(\boldsymbol{a}) = \{\boldsymbol{a}\}\)。对每个芽 \(f \in \mathcal{O}(\{\boldsymbol{a}\}; \mathrm{A})\)，有 \(\Theta_{\boldsymbol{a}}(f) = f(\boldsymbol{a})\)。

命题 3. --- 设 B 是一个交换幺 Banach 代数，\(\varphi\) 是从 A 到 B 中的连续保单位元态射。设 \(\boldsymbol{a} = (a_1, \ldots, a_n) \in \mathrm{A}^n\)。令 \(\boldsymbol{b} = (\varphi(a_1), \ldots, \varphi(a_n))\)，于是 \(\mathrm{Sp}_{\mathrm{B}}^{n}(\boldsymbol{b}) \subset \mathrm{Sp}_{\mathrm{A}}^{n}(\boldsymbol{a})\)。对每个 \(f \in \mathcal{O}(\mathrm{Sp}_{\mathrm{A}}^{n}(\boldsymbol{a}); \mathrm{A})\)，有

\[
\varphi (\Theta_ {\boldsymbol {a}} (f)) = \Theta_ {\boldsymbol {b}} (\varphi_ {*} (f)),
\]

其中 \(\varphi_{*}(f)\) 表示 \(\varphi \circ f\) 在 \(\mathrm{Sp}_{\mathrm{B}}^{n}(\boldsymbol{b})\) 的一个邻域中的芽。

只需证明：对 \(\mathrm{Sp}_{\mathrm{A}}^{n}(\boldsymbol{a})\) 的每个开邻域 U 以及每个 \(f \in \mathcal{O}(\mathrm{U}; \mathrm{A})\)，有 \(\varphi(\Theta_{\boldsymbol{a}}(f)) = \Theta_{\boldsymbol{b}}(\varphi \circ f)\)，其中 \(\varphi \circ f \in \mathcal{O}(\mathrm{U}; \mathrm{B})\)。设 \((h, u_{1}, \ldots, u_{n})\) 是适配于 \(\boldsymbol{a}\) 的序列，其中 h 的支集包含于 U (I, p.~52, 引理 1 及 I, p.~54, 引理 3)。用 \(\omega\) 表示相伴的微分形式。对每个 \(z \in C^{n}\)，有

\[
\sum_ {j = 1} ^ {n} (z _ {j} - b _ {j}) \varphi (u _ {i} (\boldsymbol {z})) = \varphi \left(\sum_ {j = 1} ^ {n} (z _ {j} - a _ {j}) u _ {j} (\boldsymbol {z})\right) = 1 - h (\boldsymbol {z}),
\]

因而序列 \((h,\varphi\circ u_{1},\ldots,\varphi\circ u_{n})\) 适配于 b。设 \(\omega'\) 是相伴的微分形式。用 \(\mu\) 表示 \(C^{n}\) 上的 Lebesgue 测度。设 \(f\in\mathcal{O}(U;A)\)。用 \(\psi\cdot\mu\) 表示与微分形式 \(f\omega\) 相伴的向量测度。与微分形式

\[
(\varphi \circ f) \omega^ {\prime} = (\varphi \circ f) d (\varphi \circ u _ {1}) \wedge d z _ {1} \wedge \dots \wedge d (\varphi \circ u _ {n}) \wedge d z _ {n}
\]

相伴的向量测度等于 \((\varphi \circ \psi) \cdot \mu\)。因此，由公式 (5) 以及 INT, VI, §2, 第 2 小节, 命题 2，我们有

\[
\Theta_ {\boldsymbol {b}} (\varphi \circ f) = \frac {n !}{(2 i \pi) ^ {n}} \int_ {\mathrm{U}} (\varphi \circ f) \mu = \frac {n !}{(2 i \pi) ^ {n}} \varphi \left(\int_ {\mathrm{U}} \psi \mu\right) = \varphi (\Theta_ {\boldsymbol {a}} (f)),
\]

这正是所要证明的。

推论 1. --- 设 \(\chi \in \mathrm{X}(\mathrm{A})\)，\(\boldsymbol{a} \in \mathrm{A}^n\)。对每个芽 \(f \in \mathcal{O}(\mathrm{Sp}^n(\boldsymbol{a}))\)，有 \(\chi(\Theta_{\boldsymbol{a}}(f)) = f(\chi(a_1), \ldots, \chi(a_n))\)。

这可由命题 3（应用于连续保单位元态射 \(\chi\colon\mathbf{A}\to\mathbf{C}\) (I, p.~29, 定理 1)）以及命题 2 的推论（应用于 Banach 代数 C）得到。

注记. --- 假设代数 A 是半单的。设 \(\boldsymbol{a} = (a_{1}, \ldots, a_{n}) \in \mathrm{A}^{n}\)。由 I, p.~38, 命题 8，映射 \(\Theta_{\boldsymbol{a}}^{\mathrm{U}}\) 是唯一的映射 \(\varphi\)：从 \(\mathcal{O}(\mathrm{U})\) 到 A 中，使得 \(\chi(\varphi(f)) = f(\chi(a_{1}), \ldots, \chi(a_{n}))\)（对一切 \(\chi \in \mathrm{X}(\mathrm{A})\) 以及每个函数 \(f \in \mathcal{O}(\mathrm{U})\)）。

推论 2. --- 设 p 是一个 ≥ 1 的整数。对每个族 \((f_{1},\ldots,f_{p})\)：由 \(\mathcal{O}(\mathrm{Sp}^{n}(\boldsymbol{a}))\) 的元素组成，有

\[
\operatorname{Sp} ^ {p} \left(\left(\Theta_ {\boldsymbol {a}} \left(f _ {1}\right), \dots , \Theta_ {\boldsymbol {a}} \left(f _ {p}\right)\right)\right) = \left(f _ {1}, \dots , f _ {p}\right) \left(\operatorname{Sp} ^ {n} (\boldsymbol {a})\right).
\]

特别地，对每个 \(f \in \mathcal{O}(\mathrm{Sp}^n(\boldsymbol{a}))\)，有 \(\mathrm{Sp}(\Theta_{\boldsymbol{a}}(f)) = f(\mathrm{Sp}^n(\boldsymbol{a}))\)。

这由推论 1 及联合谱的定义得出。

例. --- 设 A 是由单位圆上具有绝对收敛 Fourier 级数的函数所组成的复 Banach 代数 (I, p.~19, 例 8)。设 \(\varphi \in\) A。设 \(f\) 是 \(\varphi\) 的值集的某个邻域中的全纯函数芽。则 \(\psi = \Theta_{\varphi}(f)\) 是一个绝对收敛 Fourier 级数，它对每个 \(u \in \mathbf{U}\) 满足 \(\psi(u) = f(\varphi(u))\)（推论 1 应用于特征标 \(\varphi \mapsto \varphi(u)\)）。换句话说，单位圆上的函数 \(f \circ \varphi\) 也具有绝对收敛的 Fourier 级数（“P. Lévy 定理”）。这一结果推广了 Wiener 定理 (I, p.~38, 例 4)，后者涉及函数 \(f(z) = 1/z\)（在 \(\mathbf{C} - \{0\}\) 上）当 \(\varphi\) 不取零值时的情形。

